\documentclass[11pt,a4paper]{article}
\usepackage[utf8]{inputenc}
\usepackage[T1]{fontenc}
\usepackage[brazilian]{babel}
\usepackage{geometry}
\usepackage{graphicx}
\usepackage{booktabs}
\usepackage{float}
\usepackage{enumitem}
\usepackage{amsmath}
\usepackage{xcolor}
\usepackage{tcolorbox}
\geometry{margin=2.5cm}
\setlength{\parindent}{0pt}
\setlength{\parskip}{0.5em}
\tcbuselibrary{skins,breakable}
\newtcolorbox{nota}[1][]{colback=yellow!10,colframe=orange!60!black,fonttitle=\bfseries,title=#1,breakable}

\begin{document}
\begin{center}
  
\includegraphics[width=0.16\textwidth]{../../../figures_global/unb.png}\\[0.3cm]
  {\large\textbf{GABARITO DO PROFESSOR -- USO RESTRITO}}\\[0.2cm]
  {\large FGA0092 -- Princípios de Comunicação para Engenharia}\\[0.2cm]
  \rule{12cm}{0.8pt}\\[0.2cm]
  {\Large\textbf{Prática 04: Modulação SSB}}\\[0.2cm]
  \rule{12cm}{0.8pt}
\end{center}

\section*{Gráfico esperado}
\begin{center}
  \includegraphics[width=0.88\textwidth]{grafico_esperado.png}
\end{center}
\textit{O gráfico acima foi gerado pelo script \texttt{gerar\_grafico.py} desta pasta.}

\section*{Respostas esperadas}

\begin{enumerate}[label=\textbf{Q\arabic*.}, leftmargin=*, itemsep=1em]

  \item \textbf{Como identificar no espectro qual é o USB (Upper Sideband) e qual é o LSB (Lower Sideband)?}

    \textbf{Resposta esperada:} No espectro do DSB-SC com $f_c=12000$\,Hz e
    $f_m=2000$\,Hz, o USB é o pico em $f_c + f_m = 14000$\,Hz e o LSB é o pico em
    $f_c - f_m = 10000$\,Hz. A banda lateral superior (USB) está acima da portadora;
    o LSB está abaixo. Para mensagens de voz ou musicais, USB e LSB carregam as mesmas
    informações (redundantes), sendo necessário transmitir apenas uma das bandas.

    \begin{nota}[Erros comuns]
      Trocar USB e LSB, ou não identificar corretamente as frequências de corte do
      BPF usado para selecionar a banda desejada.
    \end{nota}

  \item \textbf{Qual é a diferença de largura de banda entre o DSB-SC e o SSB para a mesma mensagem?}

    \textbf{Resposta esperada:} $B_\text{DSB} = 2f_m = 4000$\,Hz (do pico em 10000 ao
    pico em 14000\,Hz). $B_\text{SSB} = f_m = 2000$\,Hz (apenas um pico em 10000 ou
    14000\,Hz). A relação $B_\text{SSB} = B_\text{DSB}/2$ é confirmada experimentalmente:
    redução exata de $50\%$ na largura de banda, que é a principal vantagem do SSB sobre
    o DSB-SC.

    \begin{nota}[Erros comuns]
      Medir a largura de banda do SSB como a distância entre os dois picos do DSB
      (4000\,Hz) ou não identificar que o SSB contém apenas um dos dois picos.
    \end{nota}

  \item \textbf{Quais são as principais vantagens e desvantagens do SSB em relação ao DSB-SC e ao AM convencional?}

    \textbf{Resposta esperada:} Vantagens: (1) ocupa metade do espectro do DSB-SC,
    permitindo dobrar o número de canais na mesma faixa; (2) maior eficiência de potência
    (toda a potência em uma banda lateral). Desvantagens: (1) demodulação requer
    oscilador local perfeitamente sincronizado; (2) geração mais complexa (BPF muito
    seletivo ou método de fase); (3) sensível a distorção de fase introduzida pelo canal.

    \begin{nota}[Erros comuns]
      Listar apenas as vantagens do SSB sem mencionar que a demodulação é mais complexa
      que no AM convencional com detector de envoltória.
    \end{nota}

  \item \textbf{Como o SSB se comporta em relação ao DSB-SC quando há erro de fase no oscilador local?}

    \textbf{Resposta esperada:} No DSB-SC, um erro de fase $\theta$ atenua o sinal
    recuperado por $\cos\theta$ sem distorção de forma de onda. No SSB, um erro de fase
    distorce a forma de onda além de atenuar, porque a transformada de Hilbert e a
    portadora devem estar em perfeita quadratura. Um erro de fase gera mistura das
    componentes I e Q, causando distorção harmônica audível em sinais de voz.

    \begin{nota}[Erros comuns]
      Afirmar que SSB e DSB-SC têm a mesma sensibilidade ao erro de fase, ignorando
      que no SSB o erro de fase causa distorção de forma de onda, não apenas atenuação.
    \end{nota}

  \item \textbf{Compare o método de filtragem com o método da transformada de Hilbert para geração de SSB.}

    \textbf{Resposta esperada:} O método de filtragem é mais simples de implementar em
    software (apenas um BPF) mas requer filtro muito seletivo quando $f_m \to 0$
    (mensagem com componentes próximas de DC). O método de Hilbert não tem essa
    restrição mas exige dois caminhos com atraso de fase exato de $-90°$ em todas as
    frequências (filtro FIR de muitos coeficientes). No GNU Radio, o método de filtragem
    é mais prático para demonstração didática.

    \begin{nota}[Erros comuns]
      Afirmar que um método é sempre superior ao outro sem contextualizar: o método
      de filtragem falha para mensagens com baixas frequências; o método de Hilbert
      tem custo computacional maior.
    \end{nota}

\end{enumerate}

\section*{Critérios de correção}
\begin{tabular}{ll}
  \toprule
  Critério & Pontuação \\
  \midrule
  Q1 -- Identificação correta de USB/LSB com frequências numéricas         & 1,0 ponto \\
  Q2 -- $B_\text{SSB} = f_m$ e $B_\text{DSB} = 2f_m$ com valores medidos  & 1,0 ponto \\
  Q3 -- Vantagens e desvantagens com menção à complexidade de demodulação  & 1,0 ponto \\
  Q4 -- Distorção de forma de onda no SSB vs.\ atenuação simples no DSB-SC & 1,0 ponto \\
  Q5 -- Comparação dos dois métodos com limitações de cada um              & 1,0 ponto \\
  \midrule
  \textbf{Total} & \textbf{5,0 pontos} \\
  \bottomrule
\end{tabular}

\bigskip
\begin{nota}[Nota de correção geral]
  Priorizar coerência física e evidência experimental (capturas de tela, tabelas).
  Respostas sem evidência visual devem ter desconto. Não penalizar pequenos erros
  de cálculo se o raciocínio físico estiver correto.
\end{nota}

\end{document}
